\paragraph{AQ-136-COSPAN: admissible path pullbacks give common sources.}
The supplied SA08 intake reports corresponding statements
A-002 and A-003, but their prior full proofs were not supplied.
The following is a complete proof in the exact category
$\mathcal A_k$ defined from the recovered original sheets.
It uses the already proved coefficient-preserving simplicial
cotensor, rather than a strict equalizer of two composites.
No novelty claim is made.

Let $k$ be any commutative ring. For every cospan of admissible
strict PD theories
\[
                         A\xrightarrow{f}C
                                      \xleftarrow{g}B
                                                               \tag{136.CS1}
\]
there is an admissible $P$ with strict maps $P\to A$ and
$P\to B$. The two resulting maps $P\to C$ are joined by
an actual simplicial cotensor homotopy. Equality of those
two maps is not asserted.

\begin{proof}
We first record the extension property used below. If a
simplicial $k$-module $F$ has acyclic normalized complex,
including $H_0=0$, then for every $n$ the restriction
\[
 \operatorname{Hom}_{s\mathrm{Mod}_k}
       (k[\Delta[n]\times\Delta[1]],F)
   \longrightarrow
 \operatorname{Hom}_{s\mathrm{Mod}_k}
       (k[\Delta[n]\times\partial\Delta[1]],F)
                                                               \tag{136.CS2}
\]
is surjective. Under Dold--Kan, the inclusion of the two
free simplicial modules normalizes to a degreewise split
inclusion of free nonnegative chain complexes. Its quotient
has a basis of relative nondegenerate simplices.
Given the chain map on the subcomplex, extend it in
increasing chain degree. For a new basis element of degree
$r\ge1$, the already assigned image of its boundary is a
cycle in $N_{r-1}F$. It is a boundary by acyclicity, so
choose its primitive in $N_rF$ as the new image. New
degree-zero basis elements, if any, can be assigned zero.
The boundary-square-zero identity verifies the cycle
condition at each step. This constructs a chain extension;
the Dold--Kan equivalence makes it the required simplicial
extension. In the displayed prism all vertices already
belong to the subcomplex. There is no dimension-dependent
condition on $F$: its one acyclicity hypothesis works for
every $n$. These linear facts follow also directly from
Stacks Project, Theorem~14.24.3, Tag~019G.

Let $P_s(C)$ be the cotensor of AQ-136-COTENSOR:
\[
                P_s(C)_n=C(\Delta[n]\times\Delta[1]).
\]
It is admissible over every $k$, with coefficient map,
endpoint maps $e_0,e_1:P_s(C)\to C$ and constant-path map
$C\to P_s(C)$ all strict. Form the actual limit of strict
PD differential algebras
\[
 \begin{gathered}
        P=(A\times B)\times_{C\times C}P_s(C),\\
 (a,b,h)\in P\quad\Longleftrightarrow\quad
                   f(a)=e_0(h),\quad g(b)=e_1(h).
 \end{gathered}                                            \tag{136.CS3}
\]
Its coefficient map sends $s$ to $(s,s,s)$ with constant
path in the third position. Limits carry the operations
componentwise, including divided powers on all positive
even elements; hence this is an actual simplicial strict
PD differential $S_k$-algebra. It has the required strict
projections and the path $h$ between their two composites.
It remains to verify admissibility with uniform cutoffs.

Fix a form degree $q$, and choose $w\ge\mu_C(q)$ so that
$F=C_*^{w,q}$ has acyclic normalization. Formula (136.CS2)
and the two disjoint copies of $\Delta[n]$ identify its
surjective restriction map with
\[
       (e_0,e_1):P_s(C)_n^{w,q}
                           \longrightarrow C_n^{w,q}\oplus C_n^{w,q}.
                                                               \tag{136.CS4}
\]
Consequently, in these weights, the defining kernel of
(136.CS3) gives an exact sequence of simplicial modules
\[
 \begin{gathered}
 0\longrightarrow P_*^{w,q}\longrightarrow
       A_*^{w,q}\oplus B_*^{w,q}\oplus P_s(C)_*^{w,q}
       \xrightarrow{\delta}
       C_*^{w,q}\oplus C_*^{w,q}\longrightarrow0,\\
       \delta(a,b,h)=(f(a)-e_0(h),\,g(b)-e_1(h)).
 \end{gathered}                                            \tag{136.CS5}
\]
The minus signs here describe a map of underlying modules;
the algebra $P$ itself was defined by a limit of algebra
maps in (136.CS3). There is no assertion that $\delta$ is
an algebra homomorphism.

For the horizontal condition, enlarge $w$ so that the
form-$q$ rows of $A,B,C,P_s(C)$ are all acyclic.
Normalization is exact, so the long exact homology sequence
of (136.CS5) proves that the row of $P$ is acyclic, including
degree zero. Its cutoff is independent of the simplicial
degree and of the horizontal homology degree. This proves
the original horizontal condition for each $q$.

For the vertical condition, fix $d\ge1$. Choose one $W_d$
so large that for every $w\ge W_d$ and every simplex $n$,
the vertical cohomology of $A_n,B_n,C_n,P_s(C)_n$ through
degree $d+1$ is canonically $S_k^w$ in degree zero and zero
in positive degrees. Such a uniform bound exists by
admissibility and the already proved cotensor comparison.
Also require $w\ge\mu_C(q)$ for all $0\le q\le d+1$.
The finite maximum of these bounds still works uniformly
in $n$.

Write $D=A_n^{w,*}\oplus B_n^{w,*}\oplus P_s(C)_n^{w,*}$,
$E=C_n^{w,*}\oplus C_n^{w,*}$, and $I=\operatorname{im}\delta$.
The sequence $0\to P_n^{w,*}\to D\to I\to0$ is exact
in all degrees, and $I\to E$ is an isomorphism in degrees
zero through $d+1$ by (136.CS4). Thus it induces an
isomorphism on cohomology through degree $d$.
In degree zero the map $H^0(D)\to H^0(I)=H^0(E)$ is
exactly
\[
             (S_k^w)^3\longrightarrow(S_k^w)^2,
                       \qquad (a,b,h)\longmapsto(a-h,b-h).
                                                               \tag{136.CS6}
\]
It is surjective and has diagonal kernel $S_k^w$.
All positive cohomology of $D$ and $I$ through degree $d$
is zero. Their long exact sequence therefore gives
\[
           H^0(P_n^{w,*})=S_k^w,\qquad
           H^q(P_n^{w,*})=0\quad(1\le q\le d).            \tag{136.CS7}
\]
The first identification is the original coefficient map.
Since the same $W_d$ works in every simplex, these are
precisely the uniform vertical tail conditions.
Together with the horizontal argument they prove
$P\in\mathcal A_k$.
\end{proof}

\paragraph{AQ-136-CONNECTED-SOURCES: the exact connectivity criterion.}
Over every commutative ring $k$, two objects of $\mathcal A_k$
admit a common admissible source if and only if they belong
to the same connected component of its underlying category,
or equivalently of its ordinary all-arrow localization.
Every finite family in one such component has a common
admissible source.

\begin{proof}
Let $A\sim B$ mean that there is an admissible $D$ with
maps $D\to A,D\to B$. This relation is reflexive and
symmetric. If $D$ is a common source of $A,B$ and $E$ a
common source of $B,C$, apply AQ-136-COSPAN to the cospan
$D\to B\leftarrow E$. Its common source $P$ maps to
$A$ through $D$ and to $C$ through $E$. Thus $\sim$ is
transitive. Every arrow $A\to B$ gives $A\sim B$, using
$A$ itself. Hence any finite zigzag gives a common source.
Conversely a common source supplies the zigzag
$A\leftarrow D\to B$. This proves the equivalence.
For a finite family, join a common source for the first
objects to the next object by this equivalence, and compose
its maps with the earlier source maps. Induction gives
one object mapping strictly to every member of the family.
The zero ring is harmless throughout.
\end{proof}

\paragraph{AQ-136-SEMILOCAL-SOURCES: strict common sources for the mixed semilocal case.}
Let $k$ be any commutative ring in which all but finitely
many rational primes are units. Every finite family in
$\mathcal A_k$ has a common admissible strict PD source.
In particular this includes $\mathbb Z_{(2,3)}$ and
every $\mathbb Z_{(P)}$-algebra for a finite set $P$ of
rational primes, with no flatness, reducedness or noetherian
hypothesis.

\begin{proof}
AQ-136-SEMILOCAL proves that the ordinary all-arrow
localization of this very category is terminal. It is
therefore connected. Apply AQ-136-CONNECTED-SOURCES.
The resulting maps to the original targets are strict
maps of admissible PD theories, even though the intermediate
cospan squares commute only by the specified cotensor
homotopies. No equality $Fh=\eta_Cu$ of the kind refuted
by AQ-136-PATH-LIFT has been used or asserted.
\end{proof}

Over $\mathbb Z$, the original common-source conjecture is
now precisely the connectedness question. The proof does
not establish that connectedness, or that every loop in
the integral all-arrow localization is trivial. Those two
integral assertions remain distinct. This result also does
not repair the prescribed partial-source maps excluded by
AQ-136-HIGHER-COMPLETION.
